\clearpage
\fancyhead[L]{\small CONNOISSANCE DES TEMS.}
\hypertarget{native-522-lower}{}\label{native-522-lower}
\begin{center}\small Göttingische gelehrte Anzeigen. No. 95. pp. 939..943. 17 June 1809.\\[3pt]\rule{.72\linewidth}{.3pt}\end{center}

\emph{Connoissance des tems ou des mouvemens célestes, à l'usage des astronomes et des navigateurs, pour l'an} 1810; \emph{publiée par le bureau des longitudes. De l'imprimerie impériale} 1808. Large octavo. Paris.

We pass over in silence the calendar and the other regular articles, whose arrangement has stayed the same. We mention only that the table of geographical longitudes and latitudes has again received some additions of no small importance, chiefly \SourceLeaf{523}several valuable new determinations from the island of Cyprus, Arabia and the Red Sea. The \emph{Additions}, which are the main reason for this review, open, as in the preceding two annual volumes, with Mr \textsc{Bouvard}’s observations at the Imperial Observatory during 1806. They fill 80 pages. Besides the Sun and Moon, we find all the older planets there, but this time no observations at all of Ceres, Pallas or Juno. \textsc{Bouvard} observed the comet \textsc{Pons} discovered in Marseille in November from 21 November to 19 December. There are fairly numerous eclipses of Jupiter’s satellites, and five stellar occultations in all, with ingresses only ($g$ in Leo on 9 January, $\zeta$ in Cancer on 1 March, $\theta$ in Ophiuchus on 1 June, $\xi$ in Gemini on 8 September, 19 in Pisces on 20 November). — Next come Chinese observations from the year 147 before our era onward, sent in by Father \textsc{Gaubil} in 1749. They are limited to close conjunctions and occultations of the Moon and planets with one another and with fixed stars, with no further details and usually with only the day given, without the hour. Little practical use can therefore be made of them. — A note by \textsc{J. C. Burckhardt} on \textsc{Hevel}’s and \textsc{Dörfel}’s contributions to the theory of cometary motion. — The sixth and final collection of \textsc{Messier}’s observations, from 1752 to the end of 1759. It includes two lunar eclipses, many occultations of stars by the Moon, still more numerous eclipses of Jupiter’s satellites, and two comets. One of these, \textsc{Halley}’s, was first found on 21 January 1759. At his teacher \textsc{Delisle}’s request, \textsc{Messier} had to keep the discovery secret until the beginning of April. — \textsc{Vidal}’s observation at Mirepoix of the superior conjunction of Mercury and the inferior conjunction of Venus in September and October 1807. — Observation of the great comet of 1807 by the same observer, from 27 November to 4 March 1808. — Observations of the same comet by \textsc{Paul Ciera} at the Royal Observatory in Lisbon (from 7 October to 29 November) and by \textsc{Olbers} in Bremen (from 8 October to 14 February 1808). — Measurement of an arc of the Earth’s meridian and of a degree perpendicular to the meridian in the East Indies, by Brigade Major \textsc{William Lambton}. This is an excerpt from the Calcutta Memoirs. The baseline measured 40006 English feet. The endpoints of the measured meridian arc were

\begin{center}\begin{tabular}{@{}lll@{}}\emph{Paudree} & Latitude & $13^\circ\,19'\,49''018$ \\ \emph{Trivandepcrum} & Latitude & $11\quad44\quad52.\,59$ \\ \end{tabular}\end{center}
\noindent From this, a degree of latitude was found to be 60495 fathoms or 56763 toises. A degree of longitude was found to be 61061 fathoms or 57294 toises. The extent of the arc from which the latter was determined is unspecified, so we cannot judge how much accuracy to assign to it. Comparing the two would give the excessively large flattening $\frac{1}{205,67}$. Compared with the French arc measurement, the degree of latitude would give a flattening of $\frac{1}{817}$. — Observations at the Lisbon Observatory by \textsc{Paul Ciera} (apart from three occultations of the star $\alpha^2$ in Cancer by the Moon, in 1807, on 20 February, 16 April and 14 October, these are all eclipses of Jupiter’s satellites). — A method for calculating transit corrections for the transit telescope, by \textsc{Delambre}. It is a very lengthy paper, though it may be useful to less experienced observers. Mr \textsc{Delambre} seems to regard the following observation as new. Comparing several transits with their right ascensions, without using the level, absolute sidereal time, or anything else related to the observer’s zenith, can bring the transit telescope into an hour circle only, and cannot bring it into the meridian. But we cannot believe that anything so obvious could ever have escaped a thoughtful astronomer. — The lunar eclipse of 4 January 1806, observed by \textsc{Vidal} at Mirepoix. — Observations at the Imperial Observatory in Marseille by \textsc{Thulis}, from 1796..\SourceLeaf{524}1806. — Observation of the great comet of 1807 at Montauban by \textsc{Duc-la Chapelle}. — Observations at Viviers by \textsc{Flaugergues} in 1807. — Mr \textsc{von Zach}’s general tables of aberration and nutation, reprinted from his well-known work. Likewise the general tables of aberration and nutation by \textsc{Gauss}, reprinted from the Monatl. Correspondenz, together with a derivation of the formulas on which they are based. — History of astronomy for 1808. This time it consists solely of excerpts from recently published astronomical works. In \textsc{von Zach}’s \emph{Tabulae speciales}, Mr \textsc{Delambre} takes the opportunity to give the formulas underlying his tables of planetary aberration, which are used again in that work. The other works are \textsc{Cagnoli}’s \emph{Star Catalogue}, \textsc{Monteiro-da-Rocha}’s \emph{Mémoires sur l'astronomie pratique}, translated from Portuguese into French, the French translation of \textsc{Cagnoli}’s trigonometry from the second edition, the first instalment of the astronomical part of \textsc{Alexander von Humboldt}’s voyage, \textsc{L. W. Pfaff}’s \emph{de tubo Culminatorio Dorpatensi brevis narratio}, \emph{Dorpati} 1808, and the \emph{Efemeridi astronomiche di Milano per l'anno} 1809. This list shows that although the volume has somewhat fewer original papers than some of its predecessors, the \emph{Connoissance des tems} is still a very valuable reference work for astronomers. In it they can expect to find news of everything in astronomy most worth knowing. — After these excerpts comes a report on the final results of the French arc measurement, extended as far as the island of Formentera. A triangle was used in which one side measured 82555 toises. This brings the whole arc from Dunkirk to Formentera to $13^\circ\,12'\,22''$ and reduces the metre from 443,296 lines to 443,2958 lines, leaving it virtually unchanged. The volume closes with some improvements to the solar and lunar tables published by the Bureau des Longitudes, a newly calculated table for the equation of time for 1800 together with the secular changes, and \textsc{Bouvard}’s meteorological observations made in 1806.
